\section{Rekapitulasi Capaian dan Peta Kompetensi}
\label{sec:rekap-capaian}

Bab ini menutup perjalanan seluruh buku. Tujuannya bukan menambah
materi baru, melainkan membantu mahasiswa melihat keseluruhan
kompetensi yang telah dibangun dan mengukur seberapa jauh
kompetensi itu dikuasai.

\subsection{Perjalanan Kompetensi}

Perjalanan belajar mata kuliah ini mengikuti urutan RPS. Bab III--VI
membangun pemahaman tentang mesin: bilangan, arsitektur BIU/EU,
register dan \eng{flag}, serta segmentasi memori. Bab VII--IX
menghubungkan program dengan mesin melalui struktur program,
instruksi transfer data, dan interupsi BIOS/DOS. Bab X--XV melatih
penulisan program: aritmatika, logika dan \eng{shift}, mode
pengalamatan, percabangan, stack/prosedur/macro, serta array dan
string. Bab XVI melatih penelusuran dan perbaikan program.

Peta kompetensi ditunjukkan pada Gambar~\ref{fig:peta-kompetensi}.

\begin{figure}[H]
    \centering
    \begin{tikzpicture}[
        cpl/.style={
            draw=obeBlue, line width=0.8pt, fill=obeBlueLight,
            rounded corners=3pt, minimum width=2.5cm,
            minimum height=0.85cm, align=center, font=\small\bfseries,
            inner sep=3pt
        },
        cpmk/.style={
            draw=obeTeal, line width=0.8pt, fill=obeTealLight,
            rounded corners=3pt, minimum width=3.1cm,
            minimum height=0.85cm, align=center, font=\small,
            inner sep=3pt
        },
        sub/.style={
            draw=obeAmber, line width=0.8pt, fill=obeAmberLight,
            rounded corners=3pt, minimum width=3.6cm,
            minimum height=0.85cm, align=center, font=\scriptsize,
            inner sep=3pt
        },
        panah/.style={-{Stealth[length=2.2mm]}, line width=0.8pt,
            obeGray}
    ]
        \node[cpl] (cpl2) at (0,3.0) {CPL02};
        \node[cpl] (cpl3) at (0,-1.5) {CPL03};

        \node[cpmk] (c1) at (4.3,4.1) {CPMK01};
        \node[cpmk] (c2) at (4.3,1.9) {CPMK02};
        \node[cpmk] (c3) at (4.3,-0.5) {CPMK03};
        \node[cpmk] (c4) at (4.3,-2.7) {CPMK04};

        \node[sub] (s1) at (9.0,5.2) {Sub-CPMK 1.1--1.4};
        \node[sub] (s2) at (9.0,2.5) {Sub-CPMK 2.1--2.3};
        \node[sub] (s3) at (9.0,-0.5) {Sub-CPMK 3.1--3.6};
        \node[sub] (s4) at (9.0,-2.7) {Sub-CPMK 4.1--4.3};

        \draw[panah] (cpl2) -- (c1);
        \draw[panah] (cpl2) -- (c2);
        \draw[panah] (cpl3) -- (c3);
        \draw[panah] (cpl3) -- (c4);
        \draw[panah] (c1) -- (s1);
        \draw[panah] (c2) -- (s2);
        \draw[panah] (c3) -- (s3);
        \draw[panah] (c4) -- (s4);
    \end{tikzpicture}
    \caption{Peta CPL, CPMK, dan kelompok Sub-CPMK pada buku ini.}
    \label{fig:peta-kompetensi}
\end{figure}

Peta ketercapaian lengkap disajikan pada
Tabel~\ref{tab:peta-ketercapaian}.

\begin{table}[H]
    \centering
    \caption{Peta ketercapaian Sub-CPMK beserta bukti belajar dan
    asesmennya.}
    \label{tab:peta-ketercapaian}
    \footnotesize
    \begin{tabular}{@{}p{1.5cm}p{4.1cm}p{4.3cm}p{3.3cm}@{}}
        \toprule
        \textbf{Sub-CPMK} & \textbf{Kompetensi} &
        \textbf{Bukti belajar} & \textbf{Asesmen} \\
        \midrule
        1.1 & Peran assembly dan sistem bilangan &
        Latihan konversi & Kuis, UTS \\
        \addlinespace
        1.2 & Arsitektur BIU dan EU &
        Diagram dan tabel \eng{trace} & UTS \\
        \addlinespace
        1.3 & Register dan \eng{flag} &
        Tabel pengamatan \eng{flag} & Kuis, UTS \\
        \addlinespace
        1.4 & Segmentasi dan alamat fisik &
        Lembar perhitungan alamat & Tugas, UTS \\
        \addlinespace
        2.1 & Struktur program dan translasi &
        Portofolio \file{.asm} $\rightarrow$ \file{.exe} &
        Tugas 1, UTS \\
        \addlinespace
        2.2 & Instruksi transfer data &
        Program transfer data & Tugas 1, UTS \\
        \addlinespace
        2.3 & Interupsi BIOS/DOS &
        Program masukan/keluaran & Kuis, UTS \\
        \addlinespace
        3.1 & Operasi aritmatika &
        Program aritmatika & Tugas 2, UAS \\
        \addlinespace
        3.2 & Operasi logika dan \eng{shift} &
        Program manipulasi bit & Kuis, UAS \\
        \addlinespace
        3.3 & Tujuh mode pengalamatan &
        Tabel alamat efektif & Tugas 3, UAS \\
        \addlinespace
        3.4 & Percabangan dan perulangan &
        Program maksimum/pengurutan & Tugas 4, UAS \\
        \addlinespace
        3.5 & Stack, prosedur, dan macro &
        Prosedur berparameter dan macro & Tugas 5, UAS \\
        \addlinespace
        3.6 & Array, string, dan pencetakan bilangan &
        Program array/string dan konversi & Tugas 6, UAS \\
        \addlinespace
        4.1 & Penelusuran per instruksi &
        Tabel \eng{trace} & Praktikum, UAS \\
        \addlinespace
        4.2 & Pemeriksaan memori dan stack &
        Laporan \eng{dump} & Praktikum, UAS \\
        \addlinespace
        4.3 & Analisis dan perbaikan \eng{hang} &
        Laporan bug dan perbaikan & Proyek, UAS \\
        \bottomrule
    \end{tabular}
\end{table}

\subsection{Bukti Belajar yang Harus Dimiliki}

Pada akhir semester, setiap mahasiswa harus dapat menunjukkan bukti
belajar berikut.

\begin{checklistbox}
\vspace{1.1em}
\begin{itemize}
    \item[\checkbox] Portofolio berisi minimal satu berkas
          \file{.asm} yang dapat dirakit untuk setiap unit materi.
    \item[\checkbox] Jurnal praktikum terisi untuk setiap sesi.
    \item[\checkbox] Enam tugas terselesaikan dan terdokumentasi.
    \item[\checkbox] Laporan \eng{debugging} memuat tabel
          \eng{trace}, temuan bug, dan bukti perbaikan.
    \item[\checkbox] Proyek akhir selesai, berjalan, dan
          dipresentasikan.
\end{itemize}
\end{checklistbox}

\begin{infobox}
\vspace{1.1em}
Bukti belajar bersifat kumulatif. Nilai akhir tidak hanya ditentukan
UTS dan UAS, tetapi juga oleh portofolio, tugas, praktikum, dan
proyek yang dipelihara sepanjang semester.
\end{infobox}
